\begin{center}
\begin{tikzpicture}[font=\small,>=stealth]
\draw[black!30,rounded corners=2pt] (0,0) rectangle (14.4,7.4);
\node[anchor=west,font=\small\bfseries] at (.5,6.8500000000000005) {一般点的两条焦距：共用同一条竖直边 $PK$};

\begin{scope}[shift={(4.2,3.6)},scale=.92]
\draw[blue,thick] (0,0) ellipse (3 and 2.4);
\draw[->,black!45] (-3.5,0)--(5.6,0) node[right] {$x$};
\draw[violet,thick] (5,-2.55)--(5,2.55) node[above] {$l:x=a^2/c$};
\coordinate (P) at (.6,2.3515101531);
\draw[teal,very thick] (-1.8,0)--(P) node[midway,above left,fill=white,inner sep=1pt] {$u$};
\draw[orange!85!black,very thick] (1.8,0)--(P) node[midway,right=3pt,fill=white,inner sep=1pt] {$v$};
\draw[dashed] (P)--(.6,0);
\draw (.6,.18)--(.78,.18)--(.78,0);
\draw[violet,dashed,thick] (P)--(5,2.3515101531);
\draw (4.82,2.3515101531)--(4.82,2.1715101531)--(5,2.1715101531);
\fill (P) circle (2pt) node[above left=5pt] {$P(x,y)$};
\fill (-1.8,0) circle (2pt) node[below=5pt] {$F_L(-c,0)$};
\fill (1.8,0) circle (2pt) node[below=5pt,fill=white,inner sep=1pt] {$F_R(c,0)$};
\fill (.6,0) circle (2pt) node[below=5pt] {$K$};
\fill (5,2.3515101531) circle (2pt) node[right=5pt] {$H$};
\end{scope}
\node[anchor=west] at (.7,.95) {$u^2=(x+c)^2+y^2$，$v^2=(x-c)^2+y^2$；相减消去同一个 $y^2$。};
\node[anchor=west] at (.7,.4) {图取椭圆上的代表点；两焦点、垂足与勾股关系同样适用于双曲线上的点。};

\end{tikzpicture}
\end{center}
